\chapter{实证分析}

本章基于完成编码的样本展开实证分析，先呈现样本互动数据和主要变量的分布情况，再通过差异检验和多元回归判断不同内容特征与传播效果之间的关系，最后通过稳健性检验考察核心结论是否稳定。

\section{描述性统计分析}

在进入具体统计结果之前，有必要先说明本章使用的数据集情况。本文实证分析所使用的数据集来自抖音平台省级文旅账号发布的 AI 文旅短视频。研究以省级文旅官方账号为主要采集对象，将标题、链接、发布时间、点赞量、评论量、收藏量、转发量、视频时长和来源省份等公开可见信息整理为基础字段，并在此基础上进一步完成 AI 介入形式、AI 介入程度、AI 技术质感、视频主体、视频主题、内容导向和视频风格等内容变量编码。经过链接去重、文旅属性筛选、AI 痕迹核验、异常样本清理和最终复核后，形成 722 条深度分析样本。该数据集既包含视频公开互动结果，也包含人工编码得到的内容特征，因此能够支持本文从“内容特征是否影响传播效果”的角度展开统计检验和回归分析。

本文最终有效样本为 722 条。样本中点赞量均值为 7249.00，中位数为 742.5；评论量均值为 1103.22，中位数为 57.5；收藏量均值为 313.71，中位数为 30.5；转发量均值为 2442.20，中位数为 42。总互动量均值为 11108.14，中位数为 907.5，均值明显高于中位数，说明互动量存在明显右偏分布，少数高互动视频获得了远高于普通视频的用户反馈。

视频时长方面，样本平均时长为 56.15 秒，中位数为 31 秒，多数 AI 文旅短视频仍然符合短视频平台的轻量化传播特征。以 2026 年 4 月 13 日为统一采集截面计算后，发布距采集日天数均值为 279.92 天，中位数为 269 天，最短为 5 天，最长为 1226 天。该结果说明样本既包含较早发布并积累较多互动的视频，也包含采集前不久发布的近期视频，因此需要在模型中控制时间累计因素。

\begin{figure}[!htbp]
  \centering
  \includegraphics[width=\textwidth]{../logs/01_log总互动量分布.png}
  \caption{log 总互动量分布}
  \label{fig:logtotal}
\end{figure}

从图 \ref{fig:logtotal} 可以看出，对总互动量取对数后，分布相较原始互动量更加平滑，有利于降低极端值对模型估计的影响。因此，本文使用 $\log(\text{总互动量}+1)$ 作为主模型因变量。

分类变量方面，AI 介入形式中，混合型最多，占 52.42\%，AI 画面占 41.49\%，AI 音频占 6.09\%。AI 介入程度中，深度介入占 47.03\%，中度介入占 43.15\%，浅度介入占 9.82\%。AI 技术质感中，普通质感占 63.49\%，精良质感占 24.34\%，粗糙质感占 12.17\%。这些分布表明，当前省级文旅 AI 短视频已较多采用画面和声音结合的混合式 AI 表达，但高质量 AI 质感仍非绝对多数。

视频内容方面，视频主体以景色为主，占 52.42\%；人和物分别占 24.62\% 和 22.96\%。视频主题中，现代城市/生活占 37.07\%，自然风光占 29.88\%，历史人文占 24.20\%。内容导向中，娱乐审美型占 74.69\%，信息实用型占 25.31\%。视频风格中，卡通/艺术化占 44.68\%，纪实/写实占 31.67\%，抽象/超现实占 23.65\%。总体来看，AI 文旅短视频明显更偏向娱乐审美和视觉表达，而不仅仅是信息服务。

从描述性统计结果还可以看到，当前省级文旅账号对于 AI 的使用已经逐渐从“尝试性应用”走向“常态化嵌入”。混合型 AI 介入和中深度 AI 介入占比较高，说明不少账号并不是只把 AI 当作单独的配音工具或局部特效工具，而是已经将其纳入整体内容生产流程。这一现象本身就说明，AI 在文旅短视频中的角色正在从边缘性技术手段转向较为稳定的生产资源。

与此同时，高质量内容的形成速度并没有完全跟上技术扩散速度。样本中普通质感和粗糙质感视频仍占较大比例，说明“会用 AI”和“把 AI 用得好”之间依然存在明显差距。对文旅账号来说，使用 AI 的门槛已经下降，但如何把 AI 转化为具有在地特色、平台传播力和视觉完成度的内容，仍然是更具挑战性的环节。也正因为如此，后续分析不能只停留在“是否使用 AI”的层面，而必须进一步进入技术质感、视频风格和时长结构等更贴近用户感知的维度。

此外，娱乐审美型内容占比明显高于信息实用型内容，这一分布特征也反映出当前文旅账号在短视频平台上的传播策略已经明显平台化。相较于传统宣传中常见的景点介绍、路线推荐和政策信息，短视频平台中的文旅内容更强调视觉吸引、情绪唤起和记忆点打造。文旅内容进入短视频平台后，不仅承担信息传播功能，也参与平台注意力竞争，这为后续回归中“风格”和“质感”变量的重要性提供了现实基础。

\section{差异检验与多元回归分析}

为初步判断不同类别之间的传播效果差异，本文对分类变量进行非参数差异检验。结果显示，所有分类变量在单因素检验中均存在显著差异。其中，AI 技术质感差异最为明显，Kruskal-Wallis H 检验的 $p$ 值为 $6.207\times 10^{-46}$，效应量为 0.2864；视频风格的 $p$ 值为 $7.313\times 10^{-32}$，效应量为 0.1964；内容导向的 $p$ 值为 $9.320\times 10^{-12}$，效应量为 0.3369。视频主题、AI 介入程度、视频主体和 AI 介入形式也均达到统计显著。

\begin{table}[htbp]
  \centering
  \caption{单因素差异检验结果}
  \label{tab:univariate}
  \begin{tabular}{lccc}
    \toprule
    变量 & 检验方法 & $p$ 值 & 效应量 \\
    \midrule
    AI 技术质感 & Kruskal-Wallis H & $6.207\times 10^{-46}$ & 0.2864 \\
    视频风格 & Kruskal-Wallis H & $7.313\times 10^{-32}$ & 0.1964 \\
    内容导向 & Mann-Whitney U & $9.320\times 10^{-12}$ & 0.3369 \\
    视频主题 & Kruskal-Wallis H & $3.678\times 10^{-8}$ & 0.0506 \\
    AI 介入程度 & Kruskal-Wallis H & $1.598\times 10^{-6}$ & 0.0343 \\
    视频主体 & Kruskal-Wallis H & $7.292\times 10^{-4}$ & 0.0173 \\
    AI 介入形式 & Kruskal-Wallis H & 0.0192 & 0.0082 \\
    \bottomrule
  \end{tabular}
\end{table}

单因素检验表明，不同 AI 特征和内容特征的视频在互动量上确实存在差异。但这一结果还不能直接证明某一变量具有独立影响，因为不同变量之间可能存在重叠关系。因此，需要进一步使用多元回归模型，在同时控制多个变量后识别各因素的净影响。

从差异强度看，AI 技术质感和视频风格在单因素检验中的表现最为突出，这说明用户面对 AI 文旅短视频时，最先产生分化反应的往往是“看起来怎么样”和“风格是否有记忆点”。相比之下，视频主题、视频主体和 AI 介入形式虽然也达到显著，但效应量相对较小，提示这些变量更像是内容表层分类，而不一定是决定传播结果的关键驱动因素。

单因素差异检验反映的是变量类别之间的表面关联，而不是控制其他因素后的净效应。例如，娱乐审美型内容在表面上可能显著高于信息实用型内容，但这种差异有可能部分来自其更常与高技术质感、强视觉风格和更短时长相结合。因此，单因素结果只能说明不同类别之间存在互动差异，仍需通过多元回归进一步判断各变量的独立影响。

本文以 $\log(\text{总互动量}+1)$ 为因变量进行主回归分析。主模型样本量为 722，自变量数为 18，$R^2$ 为 0.502，调整 $R^2$ 为 0.490。模型能够解释约 49.0\% 的总互动量对数差异。考虑到短视频互动还受到粉丝量、平台推荐、热点事件、发布时间段和投流情况等非公开因素影响，该解释力在真实平台数据研究中具有较好的有效性。

主回归结果显示，在控制其他变量后，显著影响传播效果的变量主要包括 AI 技术质感、视频风格、视频时长和发布距采集日天数。

\begin{table}[htbp]
  \centering
  \caption{主回归显著结果}
  \label{tab:mainreg}
  \begin{tabular}{lrrrr}
    \toprule
    变量 & 系数 & $p$ 值 & 影响百分比 & 方向 \\
    \midrule
    AI 技术质感：精良 & 2.857 & $2.725\times 10^{-42}$ & 1641.4\% & 正向 \\
    视频风格：抽象/超现实 & 1.543 & $3.300\times 10^{-18}$ & 368.1\% & 正向 \\
    AI 技术质感：普通 & 1.173 & $1.822\times 10^{-11}$ & 223.0\% & 正向 \\
    log 视频时长 & -0.477 & $1.285\times 10^{-9}$ & -37.9\% & 负向 \\
    log 发布距采集日天数 & 0.253 & $4.021\times 10^{-5}$ & 28.8\% & 正向 \\
    \bottomrule
  \end{tabular}
\end{table}

\begin{figure}[!htbp]
  \centering
  \includegraphics[width=\textwidth]{../logs/03_主回归显著影响百分比.png}
  \caption{主回归显著变量影响百分比}
  \label{fig:mainpercent}
\end{figure}

AI 技术质感具有最强影响。以粗糙质感为参照组，普通质感显著正向影响传播效果，影响约为 223.0\%；精良质感的影响约为 1641.4\%。这说明 AI 文旅短视频的传播效果高度依赖最终视觉质量。用户并不会因为视频使用了 AI 就自动互动，而是更容易对画面完成度高、视觉冲击强、生成痕迹较少或生成痕迹被转化为审美优势的视频产生互动。

视频风格中，抽象/超现实风格显著正向影响传播效果。以纪实/写实风格为参照组，抽象/超现实风格的影响约为 368.1\%。这一结果表明，AI 技术在文旅传播中的优势并不只是模仿真实景观，而是创造现实拍摄难以实现的视觉想象。超现实表达能够强化用户的新奇感和分享意愿，因此更容易获得平台互动。

视频时长显著负向影响传播效果。log 视频时长的影响约为 -37.9\%，意味着视频越长，互动表现越弱。这与短视频平台用户注意力有限、观看成本较高的特点有关。对于 AI 文旅内容而言，如果视觉亮点不能在较短时间内呈现，用户可能在完整观看前就划走，从而降低互动概率。

发布距采集日天数显著正向影响传播效果，影响约为 28.8\%。该变量主要反映互动累积时间。视频从发布到被采集之间经历的时间越长，越有机会获得更多点赞、评论、收藏和转发，因此在模型中控制该变量是必要的。

值得注意的是，AI 介入形式、AI 介入程度、视频主体、视频主题和内容导向在主回归中均未呈现显著影响。这并不意味着这些变量毫无意义，而是表明它们在控制技术质感、视频风格和其他因素后，没有稳定的独立解释力。

从模型结果可以进一步看出，用户可直接感知到的变量，相比内容生产过程变量更具有解释力。AI 介入形式和 AI 介入程度更接近创作端逻辑，反映的是创作者如何使用技术；AI 技术质感和视频风格则更接近接受端逻辑，反映的是用户最终看到了什么、感受到了什么。前者在控制后不显著、后者显著，说明在真实平台环境中，用户并不会因为背后用了更复杂的技术流程就自动提高互动，相反，只有当技术真正转化为更优质、更有辨识度的视觉结果时，传播优势才可能被激活。

实践中，部分明确标注 AI 的视频并没有取得较高互动。技术标签固然能够带来一定新鲜感，但如果画面粗糙、风格平淡、节奏拖沓，用户很快就会把注意力转移到别的内容上。相反，即便视频没有过度强调自身的 AI 属性，只要它在最终呈现层面足够流畅、有冲击力且有平台记忆点，就更可能获得点赞、评论与转发。AI 对文旅传播的影响主要取决于最终呈现质量，而不是技术标签本身。

发布时间变量保持显著也具有方法论意义。公开平台数据最大的局限之一，就是互动量天然受到时间累积影响。如果不控制发布时间与采集时间之间的间隔，那么早期发布的视频很容易因为累积时间更长而被误判为内容质量更高。本文将其纳入主模型，一方面是为了提高估计的严谨性，另一方面也说明时间口径统一是处理平台公开数据时不可回避的问题。

结合回归结果可以看到，本章识别出的显著变量主要对应三类影响因素。第一类是视觉完成度，即 AI 技术质感影响用户对画面协调性、自然度和专业感的评价。第二类是风格吸引力，即超现实风格通过强化视觉冲击和想象空间提高用户的停留和互动意愿。第三类是观看成本，即视频时长影响用户是否愿意继续停留并完成互动。这三类因素共同构成了本文解释 AI 文旅短视频营销效果的重要依据。

\section{高互动视频画像与地区差异分析}

本文将总互动量位于样本前 25\% 的视频定义为高互动视频，对高互动组和普通组进行比较。总互动量 P75 阈值为 4809，即总互动量达到或超过 4809 的视频被归为高互动组，共 181 条，其余样本为普通组。

高互动组和普通组在多个变量上存在显著差异。分类变量检验显示，AI 技术质感的卡方检验 $p$ 值为 $2.840\times 10^{-26}$，Cramer's V 为 0.403；视频风格的 $p$ 值为 $1.123\times 10^{-19}$，Cramer's V 为 0.347；内容导向、视频主题、AI 介入形式和 AI 介入程度也存在显著差异。

\begin{figure}[!htbp]
  \centering
  \includegraphics[width=\textwidth]{../logs/extended_work/02_高互动组画像_分类变量.png}
  \caption{高互动组与普通组的分类变量差异}
  \label{fig:highprofile}
\end{figure}

从画像结果看，高互动视频更倾向于具备较高 AI 技术质感和更强视觉风格特征。高互动组并不是简单由某一种 AI 介入形式构成，而是更突出“最终呈现质量”。这进一步说明，用户互动的关键不在于 AI 是否被明显使用，而在于 AI 是否提升了视频的审美表现和传播吸引力。

为进一步检验高互动画像结果，本文以是否进入高互动组作为因变量，构建二分类 Logit 补充模型。由于粗糙质感样本中未出现高互动视频，普通 Logit 容易出现类别分离导致系数不稳定，因此本文采用带 L2 惩罚项的 Logit 模型进行补充识别。该模型主要用于观察变量方向和相对强度，不作为替代主回归的因果检验。结果见表 \ref{tab:logit}。

\begin{table}[!htbp]
  \centering
  \caption{高互动二分类 Logit 核心结果}
  \label{tab:logit}
  \begin{tabular}{lrr}
    \toprule
    变量 & 系数 & 优势比 OR \\
    \midrule
    AI 技术质感：精良 & 2.901 & 18.194 \\
    AI 技术质感：普通 & 1.194 & 3.302 \\
    视频风格：抽象/超现实 & 1.118 & 3.058 \\
    log 视频时长 & -0.778 & 0.460 \\
    log 发布距采集日天数 & 0.530 & 1.699 \\
    \bottomrule
  \end{tabular}
\end{table}

从表 \ref{tab:logit} 可以看出，AI 技术质感和抽象/超现实风格仍然提高视频进入高互动组的可能性，视频时长则降低进入高互动组的可能性。该模型的五折交叉验证 AUC 为 0.844，说明现有内容变量对高互动视频具有一定识别能力。这一结果与前文多元回归和高互动画像基本一致，进一步表明高质感、强风格和短节奏是高互动 AI 文旅短视频的重要组合特征。

高互动视频画像进一步呈现了高互动样本的共同特征。与单独观察回归系数相比，画像分析能够更清楚地展示高互动视频在技术质感、风格表达和时长结构上的组合特点。结果表明，高互动视频往往不是在所有变量上都占优，而是在几项关键变量上形成组合优势，即更高的技术质感、更鲜明的风格表达、更紧凑的时长结构以及更强的平台记忆点。这样的组合，比单一变量更接近真实平台中的内容竞争状态。

从运营实践角度看，高互动画像也提示优质 AI 文旅短视频通常具备较强的可复用结构。例如，以地方地标或文化符号为核心视觉对象，通过拟人化、夸张化、故事化或超现实化处理，在较短时长内快速建立冲突和亮点，这类内容更容易被用户识别、记住并转发。高互动内容并不完全来自偶然流量，也可能来自较稳定的表达模式。

地区差异分析显示，不同省份/地区 AI 文旅短视频互动表现存在差异。按总互动量中位数排序，前五位分别为四川、重庆、陕西、贵州和青海。其中，四川样本量为 27，总互动量中位数为 13894，高互动占比为 51.9\%；重庆样本量为 14，总互动量中位数为 11446.5，高互动占比为 64.3\%；陕西样本量为 20，总互动量中位数为 10627，高互动占比为 70.0\%。

\begin{figure}[!htbp]
  \centering
  \includegraphics[width=\textwidth]{../logs/extended_work/01_省份互动表现Top12.png}
  \caption{省份总互动量中位数 Top 12}
  \label{fig:province}
\end{figure}

地区差异可能来自地方旅游资源知名度、账号运营能力、文旅内容创意、粉丝基础和平台推荐机制等方面。部分省份样本量较小，因此不宜过度解读其排名。本文在扩展模型中加入省份固定效应后，模型调整 $R^2$ 提升至 0.596，表明省份差异能够解释部分互动差异，但本文核心关注仍然是内容变量本身。

为更直观呈现地区差异，本文进一步绘制地区互动表现与内容特征热力矩阵。热力矩阵选取总互动量中位数排名靠前的 20 个地区，比较样本量、互动中位数、高互动占比、精良质感占比和抽象/超现实占比。由于各指标量纲不同，图中颜色深浅表示同一列指标内部的相对高低，而不是不同列之间的绝对大小。

\begin{figure}[!htbp]
  \centering
  \includegraphics[width=\textwidth]{../logs/extended_work/04_地区互动表现热力矩阵.png}
  \caption{地区互动表现与内容特征热力矩阵}
  \label{fig:provinceheat}
\end{figure}

从图 \ref{fig:provinceheat} 可以看到，互动表现较高的地区并不完全依赖样本数量，其高互动占比、精良质感占比或抽象/超现实占比往往也具有一定优势。这说明地区差异既与地方传播基础有关，也与具体内容生产方式有关。换言之，省份或地区本身能够解释一部分传播差异，但高质感和强风格仍然是影响互动表现的重要内容因素。

浙江样本具有较强典型性。该省样本数量较多，高互动案例集中出现在“AI浙江小剧场”等连续化表达中。稳定的人设化叙事和系列化运营，可能放大生成内容的传播效果。以“地标建筑热到撒腿就跑”为例，该视频总互动量达到 485737，技术质感为精良，风格为抽象/超现实，时长仅 22 秒，较好体现了高质感、强记忆点和短节奏之间的结合。

AI 文旅短视频的传播效果并不是完全脱离地方传播生态而独立存在的。不同地区在旅游资源知名度、城市话题度、账号运营能力、粉丝积累和热点承接能力上的差异，都会影响同类内容的实际传播表现。即便两条视频在视觉质量和风格上相近，也可能因为账号基础不同而获得不同的互动结果。本文通过加入省份固定效应控制这类相对稳定的地区差异，以降低地方资源优势对内容变量估计结果的干扰。

从结果来看，地区因素虽然能够提升模型解释力，但并没有改变核心内容变量的作用方向。这说明地方资源禀赋和账号基础固然会影响传播上限，但能否把这些资源转化为真正适合平台传播的 AI 内容，仍然高度依赖内容本身的生产与表达策略。对文旅账号而言，这意味着“资源好”并不自动等于“内容一定火”，内容变量仍然具有不可替代的重要性。

\section{稳健性检验}

为验证主回归结论的稳定性，本文在完成描述统计、差异检验、回归分析、高互动画像和地区差异分析之后，进一步进行稳健性检验。将稳健性检验放在本章实证分析的最后，是为了在已经识别核心影响因素之后，再判断这些结论是否会因为结果变量口径、时间累计方式或少数极端样本而发生明显变化。

本文稳健性检验并非只依赖极端值处理，而是从结果变量、时间口径、极端样本和模型设定四个方面展开，具体设计见表 \ref{tab:robustdesign}。

\begin{table}[!htbp]
  \centering
  \caption{稳健性检验设计}
  \label{tab:robustdesign}
  \begin{tabular}{L{0.22\textwidth}L{0.34\textwidth}L{0.34\textwidth}}
    \toprule
    检验方向 & 具体做法 & 检验目的 \\
    \midrule
    替换因变量 & 分别使用点赞量、评论量、收藏量和转发量的对数值重新估计模型 & 判断结论是否依赖总互动量这一综合指标 \\
    调整时间口径 & 使用日均互动量作为辅助因变量 & 判断发布时间累积效应是否改变核心结论 \\
    处理极端值 & 剔除 Top 1\%、Top 5\% 样本，并进行 1\%-99\% 缩尾处理 & 判断少数爆款视频是否主导模型结果 \\
    调整模型设定 & 加入省份固定效应 & 控制地区资源、账号基础和运营条件等相对稳定差异 \\
    \bottomrule
  \end{tabular}
\end{table}

本文分别以 $\log(\text{点赞量}+1)$、$\log(\text{评论量}+1)$、$\log(\text{收藏量}+1)$ 和 $\log(\text{转发量}+1)$ 作为因变量重复估计模型。结果显示，AI 技术质感和抽象/超现实风格在多个互动指标上仍保持显著，说明主结论并非只由某一种互动行为驱动。技术质感和风格优势不仅体现在点赞这一即时认可行为上，也会在评论、收藏和转发等不同互动形式中有所体现。

考虑到互动量是累计指标，采集时点差异可能放大不同发布时间视频之间的自然累积差距，本文进一步构造日均互动量，即总互动量除以发布距采集日天数，并以其对数形式作为替代结果变量进行辅助检验。该处理能够降低“发布时间越早、累计互动越高”对结论的机械性影响，从而更直接观察内容特征与单位时间互动表现之间的关系。结果显示，在以 $\log(\text{日均互动量}+1)$ 为因变量时，AI 技术质感中的普通和精良类别、视频风格中的抽象/超现实类别仍保持显著正向影响，视频时长仍保持显著负向影响；与此同时，发布距采集日天数转为显著负向，说明越早发布所带来的累计优势在按时间折算后被明显削弱。

在模型设定方面，本文进一步加入省份固定效应。省份固定效应用于控制不同地区在旅游资源知名度、账号基础、运营能力和地方话题度等方面的相对稳定差异。加入省份固定效应后，模型调整 $R^2$ 提升至 0.596，说明地区差异能够解释部分互动差异；但核心内容变量的解释方向并未发生明显改变，AI 技术质感和视频风格仍是理解传播效果差异的重要因素。

短视频平台中的传播数据通常存在明显长尾分布，少数爆款视频可能对模型估计产生较强影响。为排除这一问题，本文分别采用剔除总互动量 Top 1\%、剔除总互动量 Top 5\% 和对 log 总互动量进行 1\%-99\% 缩尾处理的方法重新估计模型。结果显示，AI 技术质感中的普通和精良类别，以及视频风格中的抽象/超现实类别，在不同场景下均保持显著。

\begin{figure}[!htbp]
  \centering
  \includegraphics[width=\textwidth]{../logs/extended_work/03_极端值稳健性核心系数.png}
  \caption{极端值处理后的核心变量稳健性检验}
  \label{fig:robust}
\end{figure}

在极端值处理后，模型调整 $R^2$ 分别保持在 0.459 至 0.491 之间，说明模型解释力没有因剔除爆款视频而明显失效。综合替换因变量、调整时间口径、加入省份固定效应和处理极端值的结果可以看出，AI 技术质感和抽象/超现实风格对传播效果的影响并非完全由某一种互动指标、较早发布时间、地区差异或少数爆款视频造成，而是在样本整体中具有相对稳定的解释力。

\section{本章小结}

本章围绕 722 条深度分析样本，依次进行了描述性统计、差异检验与多元回归分析、高互动视频画像与地区差异分析、稳健性检验。各部分并不是孤立展开，而是共同服务于同一个目的：先描述 AI 文旅短视频传播效果的基本分布，再识别不同内容类型之间是否存在表面差异，进一步在控制多项变量后判断哪些因素具有稳定独立影响，最后通过画像、地区比较和稳健性检验说明这些结果在具体样本和不同检验口径中的表现。

描述性统计表明，AI 文旅短视频互动效果具有明显长尾分布，少数高互动视频贡献了较多传播热度，这为后续使用对数化因变量和控制极端值提供了必要依据。单因素差异检验进一步显示，多项 AI 特征和内容特征在表面上都与互动表现有关，但这些差异可能受到变量重叠影响，因此需要通过多元回归识别净效应。主回归结果显示，在控制其他因素之后，真正稳定显著的核心因素主要是 AI 技术质感、视频风格、视频时长和发布距采集日天数，说明用户更直接回应的是最终呈现质量、风格新奇度、观看成本和互动累积时间。

高互动视频画像和地区差异分析使回归结果具有更强解释性。高互动样本往往体现出高质感、强风格和短节奏的组合优势，说明传播效果不是单一变量作用的结果，而是多种内容特征共同形成的平台竞争力。地区差异能够解释部分互动差异，但没有改变核心内容变量的方向，说明地方资源和账号基础会影响传播上限，却不能替代内容表达本身。稳健性检验进一步证明，核心结论并不依赖某一种互动指标、某一种时间口径或少数爆款样本。

因此，本章实证结果对后续讨论具有直接影响：论文的分析重点不应停留在“文旅账号是否使用 AI”，而应转向“AI 是否被转化为高质量、有风格、低观看成本且适合平台传播的内容”。这一判断也为下一章的结果讨论、典型案例分析和实践建议提供了依据，即地方文旅账号应从技术使用逻辑转向内容呈现逻辑。
